\begin{abstract}
\normalsize\noindent
% Quantifying differences across species and individuals is fundamental to many fields of biology.
% However, it remains challenging to draw detailed functional comparisons between large populations of interacting neurons.
% In particular, most analyses report average similarity scores between pairs of animals.
% For example, few analyses are able to cluster neural systems on the basis of their population activity or predict the behavioral profile of a subject from their neural statistics alone.
% Here, we introduce a general framework for comparing neural population activity in terms of \textit{shape distances}.
% This approach defines similarity in terms of explicit geometric transformations, which can be flexibly specified to obtain different measures of population-level neural similarity.
% Moreover, differences between systems are defined by a distance that is symmetric and satisfies the triangle inequality, enabling rigorous downstream analyses such as clustering and nearest-neighbor regression.
% We demonstrate this approach on datasets spanning multiple behavioral tasks (navigation, passive viewing of images, and decision making), species (mice and non-human primates) and genotypes, highlighting its potential to uncover differences across subjects and brain regions, as well as its ability to relate neural geometry to animal behavior.
%Here, we introduce a novel approach using shape metrics to compare the geometry of neural population responses. We demonstrate the application of shape metrics on various neural datasets from rodents and non-human primates, highlighting its potential to measure animal-to-animal variability, across-region representational differences, and the relationship between neural geometry and behavior. Shape metrics provide a rigorous and unifying framework, unveiling a deeper understanding of neural representational geometry beyond existing methods.


% Quantifying differences between species and individuals is fundamental to many fields of biology. In neuroscience, it remains unclear how neural population codes differ across brain regions and individuals, especially as most behavioral variables now appear decodable almost anywhere in the brain.
% Here, we leveraged shape distances, a framework that quantifies neural similarity through explicit geometric transformations. Shape distances are a proper metric, allowing an entire collection of many recordings to become a cloud of points in a common space, over which clustering and regression can be performed rigorously. Applying this approach to simulated data and diverse experimental datasets, including head direction system, visual processing across brain regions in mice and primates, and behavioral variability in a decision-making task, we find that population geometry is structured across different scales: organized by anatomy across regions within individuals, and by behavior and genotype across individuals. 

\noindent
Recent work has argued that neural codes have minimal structure, being mixed across neurons and even across brain regions. However, these conclusions are based on summaries of single neurons or single-variable decoding. Here, we investigate how the entire set of behavioral variables is jointly arranged across brain regions and individual animals using distance metrics on population-level geometry. Properties of the metric space allow us to rigorously apply techniques for clustering, regression, and low-dimensional visualization to quantify differences between regions of the brain and individual subjects.
% Here, we move beyond this by asking how the full set of behavioral variables is jointly arranged, by leveraging shape distances, which quantify neural similarity through simple geometric transformations that satisfy the axioms of a metric. Thanks to these proprieties of shape distances, we turn an arbitrarily large collection of recordings into a cloud of points in a common space, over which clustering, regression, and nearest-neighbor prediction across datasets can be performed rigorously. 
Applying this to simulated data and to recordings from dozens of individuals and regions of mice and primate brains, we find that population geometry is structured across scales. Across regions, it recovers an anatomical hierarchy, even though individual behavioral variables are decodable almost everywhere. Despite this hierarchical organization, regions did not form clear clusters and were indistinguishable from a continuum morphing of within-region geometry. Across individuals, it is organized by behavior and genotype, distinguishing animals clusters that appear identical at the single-neuron level. Comparing neural representations as geometric objects thus reveals organization that is invisible when they are compared at the level of single neurons or by decoding performance.

\end{abstract}